% ============================================================
%  CHƯƠNG 1: GIỚI THIỆU
%  Người viết: [Họ tên thành viên phụ trách]
% ============================================================

\section{Giới thiệu}
\label{sec:introduction}

\subsection{Đặt vấn đề}
Phân cụm (clustering) là một trong những bài toán cốt lõi của học máy không có
giám sát (unsupervised learning). Trong thực tế, dữ liệu quan sát thường được tạo
ra bởi nhiều cơ chế khác nhau mà ta không biết trước nhãn; mục tiêu là tự động
khám phá cấu trúc tiềm ẩn này.

% TODO: Mở rộng phần đặt vấn đề, có thể thêm ví dụ thực tế như phân đoạn ảnh,
%       nhận dạng giọng nói, phân nhóm khách hàng, …

\subsection{Mục tiêu báo cáo}
Báo cáo này có ba mục tiêu cụ thể:
\begin{enumerate}[label=(\arabic*)]
  \item Trình bày hệ thống lý thuyết về mô hình hỗn hợp xác suất, tập trung
        vào Mô hình Hỗn hợp Gaussian (GMM).
  \item Suy dẫn chi tiết thuật toán Expectation-Maximization (EM) áp dụng
        cho GMM, bao gồm chứng minh tính đảm bảo không giảm log-likelihood
        qua mỗi vòng lặp.
  \item Thực nghiệm, phân tích và so sánh với phương pháp $k$-means cổ điển.
\end{enumerate}

\subsection{Phạm vi và cấu trúc báo cáo}
Báo cáo được tổ chức như sau:
\begin{itemize}
  \item \textbf{Phần~\ref{sec:math_foundations}} – Nền tảng toán học: xác suất,
        phân phối Gaussian đa biến và bất đẳng thức Jensen.
  \item \textbf{Phần~\ref{sec:gmm}} – Định nghĩa và đặc tính của GMM.
  \item \textbf{Phần~\ref{sec:estep}–\ref{sec:mstep}} – Suy dẫn đầy đủ
        E-step và M-step.
  \item \textbf{Phần~\ref{sec:kmeans_vs_gmm}} – So sánh GMM với $k$-means.
  \item \textbf{Phần~\ref{sec:experiments}} – Thực nghiệm.
  \item \textbf{Phần~\ref{sec:conclusion}} – Kết luận và hướng phát triển.
\end{itemize}

% TODO: Nếu muốn, bổ sung một đoạn ngắn về các ứng dụng nổi bật của GMM+EM
%       (phát hiện bất thường, mô hình hóa giọng nói GMM-UBM, …)
